\section{Conclusion}
\label{sec:conclusion}

We have presented \cags{} (Class-Adaptive Guidance Strength), a training-free method that replaces the globally fixed classifier-free guidance of MGD$^3$ with per-class adaptive guidance based on cluster entropy and intra-class variance. A systematic factor ablation from an initial four-factor design identified this two-factor configuration as optimal, with mode count and inter-class separability eliminated as ineffective. We provide a theoretical explanation: under a Gaussian mixture model, the optimal per-class guidance $\lambda_c^* \propto v_c / \sigma_c^2$ is increasing in both intra-class variance and entropy, and an adaptivity gap theorem proves that the gain from per-class adaptation grows with cross-class heterogeneity---providing a foundation for the empirical scaling result from 10 to 1000 classes. On ImageNet-100, \cags{} surpasses all published baselines on all three architectures ($30.50\%$ on ResNet-18, $29.52\%$ on ResNetAP-10, $25.70\%$ on ConvNet-6), with the advantage growing to $+31$--$+44\%$ relative at 1000 classes. Cross-generator generalization to Stable Diffusion yields $+83.8\%$ on Food-101, while remaining neutral on 10-class subsets. A hybrid ablation with D4M shows no improvement over \cags{}-alone, confirming that pure-noise generation with adaptive guidance is both simpler and more effective. \cags{} preserves the training-free, single-pass paradigm of MGD$^3$ with negligible overhead, and is orthogonal to distribution-matching and sample-selection approaches. Future directions include automatic weight selection, domain-specific generators, and extension to other diffusion architectures.
